\documentclass[10pt,a4paper]{amsart}
\usepackage{amsmath,amssymb,amsthm,graphicx}
% For commutative diagrams we use Xy-pic
\usepackage[all]{xy}

% We use 2cell for 2-commutative diagrams.
\xyoption{2cell}
% Theorem environments.
%
\theoremstyle{plain}
\newtheorem{theorem}[subsection]{Theorem}
\newtheorem{proposition}[subsection]{Proposition}
\newtheorem{lemma}[subsection]{Lemma}

\theoremstyle{definition}
\newtheorem{definition}[subsection]{Definition}
\newtheorem{example}[subsection]{Example}
\newtheorem{exercise}[subsection]{Exercise}
\newtheorem{situation}[subsection]{Situation}

\theoremstyle{remark}
\newtheorem{remark}[subsection]{Remark}
\newtheorem{remarks}[subsection]{Remarks}

\numberwithin{equation}{subsection}

% Macros
%
\def\lim{\mathop{\mathrm{lim}}\nolimits}
\def\colim{\mathop{\mathrm{colim}}\nolimits}
\def\Spec{\mathop{\mathrm{Spec}}}
\def\Hom{\mathop{\mathrm{Hom}}\nolimits}
\def\Ext{\mathop{\mathrm{Ext}}\nolimits}
\def\SheafHom{\mathop{\mathcal{H}\!\mathit{om}}\nolimits}
\def\SheafExt{\mathop{\mathcal{E}\!\mathit{xt}}\nolimits}
\def\Sch{\mathit{Sch}}
\def\Mor{\mathop{\mathrm{Mor}}\nolimits}
\def\Ob{\mathop{\mathrm{Ob}}\nolimits}
\def\Sh{\mathop{\mathit{Sh}}\nolimits}
\def\NL{\mathop{N\!L}\nolimits}
\def\CH{\mathop{\mathrm{CH}}\nolimits}
\def\proetale{{pro\text{-}\acute{e}tale}}
\def\etale{{\acute{e}tale}}
\def\QCoh{\mathit{QCoh}}
\def\Ker{\mathop{\mathrm{Ker}}}
\def\Im{\mathop{\mathrm{Im}}}
\def\Coker{\mathop{\mathrm{Coker}}}
\def\Coim{\mathop{\mathrm{Coim}}}

% Boxtimes
%
\DeclareMathSymbol{\boxtimes}{\mathbin}{AMSa}{"02}

%
% Macros for moduli stacks/spaces
%
\def\QCohstack{\mathcal{QC}\!\mathit{oh}}
\def\Cohstack{\mathcal{C}\!\mathit{oh}}
\def\Spacesstack{\mathcal{S}\!\mathit{paces}}
\def\Quotfunctor{\mathrm{Quot}}
\def\Hilbfunctor{\mathrm{Hilb}}
\def\Curvesstack{\mathcal{C}\!\mathit{urves}}
\def\Polarizedstack{\mathcal{P}\!\mathit{olarized}}
\def\Complexesstack{\mathcal{C}\!\mathit{omplexes}}
% \Pic is the operator that assigns to X its picard group, usage \Pic(X)
% \Picardstack_{X/B} denotes the Picard stack of X over B
% \Picardfunctor_{X/B} denotes the Picard functor of X over B
\def\Pic{\mathop{\mathrm{Pic}}\nolimits}
\def\Picardstack{\mathcal{P}\!\mathit{ic}}
\def\Picardfunctor{\mathrm{Pic}}
\def\Deformationcategory{\mathcal{D}\!\mathit{ef}}
\newenvironment{reference}{}{}
\newenvironment{slogan}{}{}
\newenvironment{history}{}{}
\usepackage{cleveref}
\crefname{originallemma}{original lemma}{original lemmas}
\Crefname{originallemma}{Original lemma}{Original lemmas}
\makeatletter
\newlabel{decent-spaces-lemma-bounded-fibres}{{03JX}{source}{}{sourcepointer-1}{}}
\newlabel{decent-spaces-lemma-bounded-fibres@cref}{{[originallemma][0][]03JX}{[1][0][]source}{}{}{}}
\makeatother
\crefname{originallemma}{original lemma}{original lemmas}
\Crefname{originallemma}{Original lemma}{Original lemmas}
\makeatletter
\newlabel{decent-spaces-lemma-quasi-compact-reasonable-stratification}{{07S8}{source}{}{sourcepointer-2}{}}
\newlabel{decent-spaces-lemma-quasi-compact-reasonable-stratification@cref}{{[originallemma][0][]07S8}{[1][0][]source}{}{}{}}
\makeatother
\crefname{originallemma}{original lemma}{original lemmas}
\Crefname{originallemma}{Original lemma}{Original lemmas}
\makeatletter
\newlabel{spaces-morphisms-lemma-etale-unramified}{{06CR}{source}{}{sourcepointer-3}{}}
\newlabel{spaces-morphisms-lemma-etale-unramified@cref}{{[originallemma][0][]06CR}{[1][0][]source}{}{}{}}
\makeatother
\crefname{originallemma}{original lemma}{original lemmas}
\Crefname{originallemma}{Original lemma}{Original lemmas}
\makeatletter
\newlabel{spaces-morphisms-lemma-diagonal-unramified-morphism}{{05W1}{source}{}{sourcepointer-4}{}}
\newlabel{spaces-morphisms-lemma-diagonal-unramified-morphism@cref}{{[originallemma][0][]05W1}{[1][0][]source}{}{}{}}
\makeatother
\crefname{originallemma}{original lemma}{original lemmas}
\Crefname{originallemma}{Original lemma}{Original lemmas}
\makeatletter
\newlabel{morphisms-lemma-quasi-affine-finite-type-over-S}{{04II}{source}{}{sourcepointer-5}{}}
\newlabel{morphisms-lemma-quasi-affine-finite-type-over-S@cref}{{[originallemma][0][]04II}{[1][0][]source}{}{}{}}
\makeatother
\crefname{originallemma}{original lemma}{original lemmas}
\Crefname{originallemma}{Original lemma}{Original lemmas}
\makeatletter
\newlabel{morphisms-lemma-H-quasi-projective-open-H-projective}{{01WA}{source}{}{sourcepointer-6}{}}
\newlabel{morphisms-lemma-H-quasi-projective-open-H-projective@cref}{{[originallemma][0][]01WA}{[1][0][]source}{}{}{}}
\makeatother
\crefname{originallemma}{original lemma}{original lemmas}
\Crefname{originallemma}{Original lemma}{Original lemmas}
\makeatletter
\newlabel{morphisms-lemma-quasi-compact-immersion}{{01RG}{source}{}{sourcepointer-7}{}}
\newlabel{morphisms-lemma-quasi-compact-immersion@cref}{{[originallemma][0][]01RG}{[1][0][]source}{}{}{}}
\makeatother
\crefname{originalsection}{original section}{original sections}
\Crefname{originalsection}{Original section}{Original sections}
\makeatletter
\newlabel{morphisms-section-projective}{{01W7}{source}{}{sourcepointer-8}{}}
\newlabel{morphisms-section-projective@cref}{{[originalsection][0][]01W7}{[1][0][]source}{}{}{}}
\makeatother
\crefname{originallemma}{original lemma}{original lemmas}
\Crefname{originallemma}{Original lemma}{Original lemmas}
\makeatletter
\newlabel{spaces-morphisms-lemma-equality-of-morphisms}{{084N}{source}{}{sourcepointer-9}{}}
\newlabel{spaces-morphisms-lemma-equality-of-morphisms@cref}{{[originallemma][0][]084N}{[1][0][]source}{}{}{}}
\makeatother
\crefname{originallemma}{original lemma}{original lemmas}
\Crefname{originallemma}{Original lemma}{Original lemmas}
\makeatletter
\newlabel{spaces-morphisms-lemma-refined-valuative-criterion-universally-closed}{{089G}{source}{}{sourcepointer-10}{}}
\newlabel{spaces-morphisms-lemma-refined-valuative-criterion-universally-closed@cref}{{[originallemma][0][]089G}{[1][0][]source}{}{}{}}
\makeatother
\crefname{originallemma}{original lemma}{original lemmas}
\Crefname{originallemma}{Original lemma}{Original lemmas}
\makeatletter
\newlabel{spaces-morphisms-lemma-lift-valuation-ring-through-flat-morphism}{{089F}{source}{}{sourcepointer-11}{}}
\newlabel{spaces-morphisms-lemma-lift-valuation-ring-through-flat-morphism@cref}{{[originallemma][0][]089F}{[1][0][]source}{}{}{}}
\makeatother
\crefname{originallemma}{original lemma}{original lemmas}
\Crefname{originallemma}{Original lemma}{Original lemmas}
\makeatletter
\newlabel{morphisms-lemma-locally-projective-proper}{{01WC}{source}{}{sourcepointer-12}{}}
\newlabel{morphisms-lemma-locally-projective-proper@cref}{{[originallemma][0][]01WC}{[1][0][]source}{}{}{}}
\makeatother
\crefname{originalproposition}{original proposition}{original propositions}
\Crefname{originalproposition}{Original proposition}{Original propositions}
\makeatletter
\newlabel{schemes-proposition-characterize-universally-closed}{{01KF}{source}{}{sourcepointer-13}{}}
\newlabel{schemes-proposition-characterize-universally-closed@cref}{{[originalproposition][0][]01KF}{[1][0][]source}{}{}{}}
\makeatother
\title{A proper surjective quasi-projective scheme cover of a separated finite-type algebraic space}
\author{The Stacks Project Authors}
\date{stacks-project:spaces-cohomology.tex:local-snapshot-20261011; GNU FDL 1.2 or later}
\begin{document}
\maketitle
\noindent\textit{Editorial note.} Public016 is Chow for separated finite-type schemes over a Noetherian base; public492 is a derived local modification of a finitely presented scheme. Here the source may be an algebraic space and the base ring is arbitrary; the constructed cover is an actual scheme.\par\medskip
\noindent\textit{Editorial note.} The older two substantive reviews support the complete109-line author proof; they are reused without repeating a full proof audit.\par\medskip
\noindent\textit{Allowed standard prerequisites (editorial).}\begin{itemize}
\item Explicit prerequisite boundary. Algebraic spaces and etale presentations, bounded fibre cardinality and stratification, scheme-theoretic images, projective compactification of affine schemes, equality of maps from schematically dense opens, and the refined valuative criterion.
\end{itemize}\medskip
\ifdefined\hypertarget\hypertarget{sourcepointer-1}{}\fi
\noindent\textit{Source reference.} Original lemma 03JX: focused-seeds/AGP/sources/stacks-snapshot-20261011/tags.txt, source lines 4655--4655. Stable Stacks tag for an expressly allowed foundational or notational input; the counted construction remains the two concrete core obligations. The retained local core is not replaced by this pointer. See standard\_prerequisites.\par\medskip
\ifdefined\hypertarget\hypertarget{sourcepointer-2}{}\fi
\noindent\textit{Source reference.} Original lemma 07S8: focused-seeds/AGP/sources/stacks-snapshot-20261011/tags.txt, source lines 9893--9893. Stable Stacks tag for an expressly allowed foundational or notational input; the counted construction remains the two concrete core obligations. The retained local core is not replaced by this pointer. See standard\_prerequisites.\par\medskip
\ifdefined\hypertarget\hypertarget{sourcepointer-3}{}\fi
\noindent\textit{Source reference.} Original lemma 06CR: focused-seeds/AGP/sources/stacks-snapshot-20261011/tags.txt, source lines 8155--8155. Stable Stacks tag for an expressly allowed foundational or notational input; the counted construction remains the two concrete core obligations. The retained local core is not replaced by this pointer. See standard\_prerequisites.\par\medskip
\ifdefined\hypertarget\hypertarget{sourcepointer-4}{}\fi
\noindent\textit{Source reference.} Original lemma 05W1: focused-seeds/AGP/sources/stacks-snapshot-20261011/tags.txt, source lines 7570--7570. Stable Stacks tag for an expressly allowed foundational or notational input; the counted construction remains the two concrete core obligations. The retained local core is not replaced by this pointer. See standard\_prerequisites.\par\medskip
\ifdefined\hypertarget\hypertarget{sourcepointer-5}{}\fi
\noindent\textit{Source reference.} Original lemma 04II: focused-seeds/AGP/sources/stacks-snapshot-20261011/tags.txt, source lines 5904--5904. Stable Stacks tag for an expressly allowed foundational or notational input; the counted construction remains the two concrete core obligations. The retained local core is not replaced by this pointer. See standard\_prerequisites.\par\medskip
\ifdefined\hypertarget\hypertarget{sourcepointer-6}{}\fi
\noindent\textit{Source reference.} Original lemma 01WA: focused-seeds/AGP/sources/stacks-snapshot-20261011/tags.txt, source lines 2475--2475. Stable Stacks tag for an expressly allowed foundational or notational input; the counted construction remains the two concrete core obligations. The retained local core is not replaced by this pointer. See standard\_prerequisites.\par\medskip
\ifdefined\hypertarget\hypertarget{sourcepointer-7}{}\fi
\noindent\textit{Source reference.} Original lemma 01RG: focused-seeds/AGP/sources/stacks-snapshot-20261011/tags.txt, source lines 2301--2301. Stable Stacks tag for an expressly allowed foundational or notational input; the counted construction remains the two concrete core obligations. The retained local core is not replaced by this pointer. See standard\_prerequisites.\par\medskip
\ifdefined\hypertarget\hypertarget{sourcepointer-8}{}\fi
\noindent\textit{Source reference.} Original section 01W7: focused-seeds/AGP/sources/stacks-snapshot-20261011/tags.txt, source lines 2472--2472. Stable Stacks tag for an expressly allowed foundational or notational input; the counted construction remains the two concrete core obligations. The retained local core is not replaced by this pointer. See standard\_prerequisites.\par\medskip
\ifdefined\hypertarget\hypertarget{sourcepointer-9}{}\fi
\noindent\textit{Source reference.} Original lemma 084N: focused-seeds/AGP/sources/stacks-snapshot-20261011/tags.txt, source lines 10328--10328. Stable Stacks tag for an expressly allowed foundational or notational input; the counted construction remains the two concrete core obligations. The retained local core is not replaced by this pointer. See standard\_prerequisites.\par\medskip
\ifdefined\hypertarget\hypertarget{sourcepointer-10}{}\fi
\noindent\textit{Source reference.} Original lemma 089G: focused-seeds/AGP/sources/stacks-snapshot-20261011/tags.txt, source lines 10496--10496. Stable Stacks tag for an expressly allowed foundational or notational input; the counted construction remains the two concrete core obligations. The retained local core is not replaced by this pointer. See standard\_prerequisites.\par\medskip
\ifdefined\hypertarget\hypertarget{sourcepointer-11}{}\fi
\noindent\textit{Source reference.} Original lemma 089F: focused-seeds/AGP/sources/stacks-snapshot-20261011/tags.txt, source lines 10495--10495. Stable Stacks tag for an expressly allowed foundational or notational input; the counted construction remains the two concrete core obligations. The retained local core is not replaced by this pointer. See standard\_prerequisites.\par\medskip
\ifdefined\hypertarget\hypertarget{sourcepointer-12}{}\fi
\noindent\textit{Source reference.} Original lemma 01WC: focused-seeds/AGP/sources/stacks-snapshot-20261011/tags.txt, source lines 2477--2477. Stable Stacks tag for an expressly allowed foundational or notational input; the counted construction remains the two concrete core obligations. The retained local core is not replaced by this pointer. See standard\_prerequisites.\par\medskip
\ifdefined\hypertarget\hypertarget{sourcepointer-13}{}\fi
\noindent\textit{Source reference.} Original proposition 01KF: focused-seeds/AGP/sources/stacks-snapshot-20261011/tags.txt, source lines 2048--2048. Stable Stacks tag for an expressly allowed foundational or notational input; the counted construction remains the two concrete core obligations. The retained local core is not replaced by this pointer. See standard\_prerequisites.\par\medskip
\begin{lemma}
\label{lemma-weak-chow}
Let $A$ be a ring. Let $X$ be an algebraic space over $\Spec(A)$
whose structure morphism $X \to \Spec(A)$ is separated of finite type.
Then there exists a proper surjective morphism $X' \to X$
where $X'$ is a scheme which is H-quasi-projective over $\Spec(A)$.
\end{lemma}

\begin{proof}
Let $W$ be an affine scheme and let $f : W \to X$ be a surjective
\'etale morphism. There exists an integer $d$ such that all geometric
fibres of f have $\leq d$ points (because $X$ is a separated algebraic
hence reasonable, see
Decent Spaces, Lemma \ref{decent-spaces-lemma-bounded-fibres}).
Picking $d$ minimal we get a nonempty open $U \subset X$ such that
$f^{-1}(U) \to U$ is finite \'etale of degree $d$, see
Decent Spaces, Lemma
\ref{decent-spaces-lemma-quasi-compact-reasonable-stratification}.
Let
$$
V \subset W \times_X W \times_X \ldots \times_X W
$$
($d$ factors in the fibre product) be the complement of all the diagonals.
Because $W \to X$ is separated the diagonal $W \to W \times_X W$ is a
closed immersion. Since $W \to X$ is \'etale the diagonal
$W \to W \times_X W$ is an open immersion, see
Morphisms of Spaces, Lemmas
\ref{spaces-morphisms-lemma-etale-unramified} and
\ref{spaces-morphisms-lemma-diagonal-unramified-morphism}.
Hence the diagonals are open and closed subschemes
of the quasi-compact scheme $W \times_X \ldots \times_X W$.
In particular we conclude $V$ is a quasi-compact scheme.
Choose an open immersion $W \subset Y$ with $Y$ H-projective over
$A$ (this is possible as $W$ is affine and of finite type over $A$;
for example we can use
Morphisms, Lemmas
\ref{morphisms-lemma-quasi-affine-finite-type-over-S} and
\ref{morphisms-lemma-H-quasi-projective-open-H-projective}).
Let
$$
Z \subset Y \times_A Y \times_A \ldots \times_A Y
$$
be the scheme theoretic image of the composition
$V \to W \times_X \ldots \times_X W \to Y \times_A \ldots \times_A Y$.
Observe that this morphism is quasi-compact since $V$ is quasi-compact
and $Y \times_A \ldots \times_A Y$ is separated.
Note that $V \to Z$ is an open immersion as
$V \to Y \times_A \ldots \times_A Y$ is an immersion, see
Morphisms, Lemma \ref{morphisms-lemma-quasi-compact-immersion}.
The projection morphisms give $d$ morphisms $g_i : Z \to Y$.
These morphisms $g_i$ are projective as $Y$ is projective over $A$, see
material in Morphisms, Section \ref{morphisms-section-projective}.
We set
$$
X' = \bigcup g_i^{-1}(W) \subset Z
$$
There is a morphism $X' \to X$ whose restriction to $g_i^{-1}(W)$ is
the composition $g_i^{-1}(W) \to W \to X$.
Namely, these morphisms agree over $V$ hence agree over
$g_i^{-1}(W) \cap g_j^{-1}(W)$ by
Morphisms of Spaces, Lemma \ref{spaces-morphisms-lemma-equality-of-morphisms}.
Claim: the morphism $X' \to X$ is proper.

\medskip\noindent
If the claim holds, then the lemma follows by induction on $d$.
Namely, by construction $X'$ is H-quasi-projective over $\Spec(A)$.
The image of $X' \to X$ contains the open $U$ as $V$ surjects onto $U$.
Denote $T$ the reduced induced algebraic space structure on $X \setminus U$.
Then $T \times_X W$ is a closed subscheme of $W$, hence affine.
Moreover, the morphism $T \times_X W \to T$ is \'etale and every geometric
fibre has $< d$ points. By induction hypothesis there exists a proper
surjective morphism $T' \to T$ where $T'$ is a scheme H-quasi-projective
over $\Spec(A)$. Since $T$ is a closed subspace of $X$ we see that
$T' \to X$ is a proper morphism. Thus the lemma follows by taking the
proper surjective morphism $X' \amalg T' \to X$.

\medskip\noindent
Proof of the claim. By construction the morphism $X' \to X$ is separated
and of finite type. We will check conditions (1) -- (4) of
Morphisms of Spaces, Lemma
\ref{spaces-morphisms-lemma-refined-valuative-criterion-universally-closed}
for the morphisms $V \to X'$ and $X' \to X$.
Conditions (1) and (2) we have seen above.
Condition (3) holds as $X' \to X$ is separated (as a morphism whose
source is a separated algebraic space). Thus it suffices to check
liftability to $X'$ for diagrams
$$
\xymatrix{
\Spec(K) \ar[r] \ar[d] & V \ar[d] \\
\Spec(R) \ar[r] & X
}
$$
where $R$ is a valuation ring with fraction field $K$.
Note that the top horizontal map is given by $d$ pairwise distinct
$K$-valued points $w_1, \ldots, w_d$ of $W$. In fact, this
is a complete set of inverse images of the point $x \in X(K)$
coming from the diagram. Since $W \to X$ is surjective,
we can, after possibly replacing $R$ by an extension of valuation rings,
lift the morphism $\Spec(R) \to X$ to a morphism $w : \Spec(R) \to W$, see
Morphisms of Spaces, Lemma
\ref{spaces-morphisms-lemma-lift-valuation-ring-through-flat-morphism}.
Since $w_1, \ldots, w_d$ is a complete collection of inverse images of
$x$ we see that $w|_{\Spec(K)}$ is equal to one of them, say $w_i$.
Thus we see that we get a commutative diagram
$$
\xymatrix{
\Spec(K) \ar[r] \ar[d] &  Z \ar[d]_{g_i}\\
\Spec(R) \ar[r]^w & Y
}
$$
By the valuative criterion of properness for the projective
morphism $g_i$ we can lift $w$ to $z : \Spec(R) \to Z$, see
Morphisms, Lemma \ref{morphisms-lemma-locally-projective-proper}
and
Schemes, Proposition \ref{schemes-proposition-characterize-universally-closed}.
The image of $z$ is in $g_i^{-1}(W) \subset X'$ and the proof is complete.
\end{proof}

\bibliographystyle{plain}
\bibliography{my}
\end{document}
